% !PRINTED 99
\weberpage{99}
%<<CONT
\weq{22}{4m = (2a - b)^2 + 3b^2,}

und wenn $m$ durch 9 theilbar ist, so ist $2a - b$ nach (19) durch 3
theilbar, und aus (22) folgt dann, dass auch in diesem Falle $b$ durch 3
theilbar ist. Wir können also in allen Fällen setzen

\[ 2a - b = A, \quad b = 3B, \]

so dass $A, B$ ganze Zahlen sind, die der Bedingung

\weq{23}{4m = A^2 + 27B^2}

genügen, und die Gleichung für $\xi$ wird dann

\weq{24}{\xi^3 - \frac{m}{3}\xi - \frac{mA}{27} = 0,}

und für die Discriminante dieser Gleichung folgt

\weq{25}{\sqrt{D} = mB.}

Um für die einfachsten Beispiele die Rechnung durchzuführen, nehmen wir für
$q_1 = 9, q_2 = 7, q_3 = 13$ aus \S.\,172 des ersten Bandes die Werthe:

\[ \begin{aligned}
q_1 &= 9, & a + b\varrho &= 3\varrho, \\
q_2 &= 7, & a + b\varrho &= -(1 + 3\varrho), \\
q_3 &= 13, & a + b\varrho &= -(4 + 3\varrho).
\end{aligned} \]

Daraus ergeben sich für $m = 63$ die beiden Werthe:

\[ \begin{aligned}
m = 63:\quad a + b\varrho &= -3\varrho(1 + 3\varrho) = 9 + 6\varrho, \\
a + b\varrho &= -3\varrho(1 + 3\varrho^2) = -9 - 3\varrho,
\end{aligned} \]

für

\[ \begin{aligned}
m = 91:\quad a + b\varrho &= (1 + 3\varrho)(4 + 3\varrho) = -5 + 6\varrho, \\
a + b\varrho &= (1 + 3\varrho)(4 + 3\varrho^2) = 10 + 9\varrho,
\end{aligned} \]

und für $m = 819$ die Werthe:

\[ \begin{aligned}
a + b\varrho &= 3\varrho(1 + 3\varrho)(4 + 3\varrho) = -3(6 + 11\varrho), \\
&= 3\varrho(1 + 3\varrho)(4 + 3\varrho^2) = -3(9 - \varrho), \\
&= 3\varrho(1 + 3\varrho^2)(4 + 3\varrho) = 3(9 + 10\varrho), \\
&= 3\varrho(1 + 3\varrho^2)(4 + 3\varrho^2) = 3(6 - 5\varrho).
\end{aligned} \]

Daraus finden sich die cubischen Gleichungen für $\eta$ in diesen drei
Fällen:

\[ \begin{aligned}
m = 63:\quad &\eta^3 - 21\eta - 28 = 0, & (\eta = \xi) \\
&\eta^3 - 21\eta + 35 = 0, \\
m = 91:\quad &\eta^3 - \eta^2 - 30\eta + 64 = 0, & (\eta = \xi + \tfrac{1}{3}) \\
&\eta^3 - \eta^2 - 30\eta - 27 = 0,
\end{aligned} \]
%CONT>>
